\scanpage{117}
Je vais donner une expression de $\rho$. D'abord en termes d'une base
\[
1,u,v\qquad\text{base de } B=H+H'
\]
posant
\[
\begin{aligned}
&\left\{\begin{aligned} b&=\operatorname{Tr}u,\ c=\det u\\ b'&=\operatorname{Tr}v,\ c'=\det v\\ t&=\operatorname{Tr}uv=\operatorname{Tr}vu
\end{aligned}\right. &&
\begin{aligned}\delta_H&=b^2-4c\\ \delta_{H'}&=b'^2-4c'\\ &\end{aligned}
\end{aligned}
\]
on trouve l'expression \margin{\uline{NB} la condition $H$ étale $\Longleftrightarrow\delta_H$ inv $H'$ étale $\Longleftrightarrow\delta_{H'}$ inv.}
\[
\boxed{\rho_0=\frac{P\deleted{\unsure{\delta|6}}}{Q}}
\]
où
\[
\begin{aligned}
P&=P(b,c,b',c',t)=(b^2-4c)(b'^2-4c')=\delta_H\delta_{H'}^{\unsure{*|\cdot}}\\ Q&=Q(b,c,b',c',t)=t^2-bb't+(b^2c'+b'^2c-4cc')\deleted{\unsure{\delta_{\!B}}}
\end{aligned}
\]
et d'autre part on a
\[
\begin{aligned}
Q&=\unsure{-|{\deleted{-}}}\tfrac12\delta_{\varphi|B}\,\substack{\text{pour la base}\\1,u,v} =\unsure{-|{\deleted{-}}}\tfrac12\det\begin{pmatrix} \operatorname{Tr}1 & \operatorname{Tr}u & \operatorname{Tr}v\\ \operatorname{Tr}u & \operatorname{Tr}u^2 & \operatorname{Tr}uv\\ \operatorname{Tr}v & \operatorname{Tr}vu & \operatorname{Tr}v^2
\end{pmatrix}\\
&=\unsure{-|{\deleted{-}}}\,\text{discr. divisé de la forme}\\ &\qquad\det_{|B}\ \text{p.r. à la base }1,u,v
\end{aligned}
\]
(matrice \inserted{de $\varphi_{|B}$} égale :
\[
\begin{pmatrix}2&b&b'\\ b&\deleted{\unsure{2}}b^2-2c&t\\ b'&t&b'^2-2c'\end{pmatrix}\qquad)
\]
donc par hypo ($H,H'$ en position générale) \deleted{donc} on a bien $Q$ inversible, donc $\rho_0$ est défini, et non nul.

Je voudrais donner à cette formule pour $\rho_0$ une forme intrinsèque, \unsure{indépendant|indépendante} du choix des bases.
% archive sheet number: 116
